\section{Адаптивная остановка и полный учёт запросов}
\label{nloracle:section}
Сертификаты выше совместимы с обычным квадратичным поиском шага.
Докажем общую оценку работы для калибровки $\Phi$.
Значения функции считаются точными.
Пусть $D\ge\operatorname{diam}Q>0$ и
\begin{equation}\label{nloracle:model}
 \lambda_{\max}(\nabla^2f(x))\le K(\delta(x)),\qquad
 \delta(x)\le\Phi(G(x)).
\end{equation}
Функция $K:[0,\infty)\to[0,\infty)$ конечна, непрерывна и не убывает,
а $\Phi:[0,\infty)\to[0,\infty)$ непрерывна, строго возрастает,
не ограничена сверху и удовлетворяет $\Phi(0)=0$.
Точной огибающей зазора соответствуют $\Phi=\mathcal B$ и
$K(u)=H_0+H_1u$. Условия \eqref{nloracle:model} допускают
и другие классы функций.
При известной границе диаметра можно также взять $\Phi=\mathcal B_D$.

При запрошенной точности направлений $\eta>0$ оракул возвращает
$h\ge0$ и $d=s-x$, $s\in Q$, для которых
\begin{equation}\label{nloracle:query}
 G(x)\le B=h+\eta,\qquad
 -\langle\nabla f(x),d\rangle\ge h-\eta.
\end{equation}
Сертификат в текущей точке должен удовлетворять
$\delta(x)\le C\le\Phi(B)$. Можно взять $C=\Phi(B)$
или более точную совместную границу.
Обновим $\ell\leftarrow\max\{\ell,f(x)-C\}$ и положим $U=f(x)-\ell$.
Первый запрос создаёт начальную нижнюю границу.

\paragraph{Правило.}
Выберем $\eta_0,M>0$, $0<\theta<1$ и будем переносить $\eta$
между точками. После каждого запроса сначала остановимся, если
$U\le\varepsilon$, где $\varepsilon>0$.
Иначе проверим $\eta\le\theta h$. Если условие не выполнено,
уменьшим $\eta$ вдвое и повторим запрос в той же точке,
используя сохранённое значение функции.
После успешной проверки положим $a=h-\eta$,
начнём с $L=L_{k-1}/2$, $L_{-1}=M$, и будем удваивать $L$, пока
\begin{equation}\label{eq:btstep}
 \alpha=\min\{1,a/(L\|d\|^2)\},\qquad
 f(x+\alpha d)\le f(x)-\alpha a+\tfrac12L\alpha^2\|d\|^2.
\end{equation}
Примем точку $x+\alpha d$, сохраним её значение и принятое $L_k$.
Сам поиск шага стандартен~\cite{pedregosa}; сертификаты определяют,
когда полученная точка достаточно точна.

\begin{theorem}[Остановка с переносом точности]\label{nloracle:main}
Каждый цикл уточнения и поиска шага конечен. Правило останавливается
с $\delta(x_N)\le U_N\le\varepsilon$.
Пусть $N$ есть число принятых обновлений, а $R$ есть число делений
точности пополам, включая начальную и последнюю проверки.
При $\kappa=(1-\theta)/(1+\theta)$
\begin{equation}\label{nloracle:refinement}
 R\le\left\lceil\log_2\max\left\{1,
 \frac{(1+\theta)\eta_0}{\theta\Phi^{-1}(\varepsilon)}\right\}\right\rceil,
 \qquad T=N+1+R.
\end{equation}
Запросов в текущих точках ровно $T$.
Каждый векторный запрос с полным LMO требует одного вызова LMO;
каждый запрос оценок атомов ниже требует одного полного прохода по ним.

Если нужно хотя бы одно обновление, пусть $U_0>\varepsilon$ есть
сертификат после начальной проверки точности. Положим
\[
 \begin{gathered}
 J=\lceil\log_2(U_0/\varepsilon)\rceil,\quad b_j=U_0/2^j,\quad
 s_j=\kappa\Phi^{-1}(b_j/2),\\
 C_j=\max\{2K(b_j),s_j/D^2\},\quad
 P_K=\left\lceil\log_2\frac{\max\{M,C_0\}}{C_{J-1}}\right\rceil,
 \qquad 0\le j<J.
 \end{gathered}
\]
Тогда
\begin{equation}\label{nloracle:work}
 N\le T_\varepsilon:=P_K+
 \left\lceil\sum_{j=0}^{J-1}
 \max\left\{\frac{2b_j}{s_j},\frac{4D^2K(b_j)b_j}{s_j^2}\right\}\right\rceil.
\end{equation}
При фактическом $N\ge1$ число вычислений функции, включая
начальное значение и отклонённые пробы, удовлетворяет
\begin{equation}\label{nloracle:functions}
 N_f\le1+2N+
 \left\lceil\log_2\max\{1,2K(U_0)/M\}\right\rceil.
\end{equation}
При $N=0$ требуется одно значение функции и $1+R$ запросов.
Фазовые величины в этом случае не нужны.
\end{theorem}

\begin{proof}
Каждая сохранённая нижняя граница верна, поэтому
$0\le\delta\le U\le C\le\Phi(B)$.
Если остановка ещё не выполнена, а проверка точности не прошла, то
\[
 \Phi^{-1}(\varepsilon)<\Phi^{-1}(U)\le B
   <(1+\theta)\eta/\theta.
\]
Точность уменьшается вдвое и не сбрасывается. Отсюда следуют
граница для $R$ и конечность уточнений.
В каждой посещённой точке один запрос завершает проверку,
а каждому делению точности соответствует один неудачный запрос.
Значит, $T=N+1+R$.
После успешной проверки перед обновлением
\begin{equation}\label{compact:certified-slope}
 a=h-\eta\ge\kappa(h+\eta)\ge\kappa\Phi^{-1}(U)>0.
\end{equation}

Для поиска шага положим $w(t)=f(x+td)-f(x)$,
$K_x=K(\delta(x))$.
Пока $w\le0$, имеем $w''\le K_x\|d\|^2$ и $w'(0)\le-a<0$.
Первый возврат к нулю до $2a/(K_x\|d\|^2)$ противоречил бы
$w(t)\le-at+K_x\|d\|^2t^2/2<0$.
Поэтому любое $L\ge K_x$ проходит проверку с предложенным шагом.
При $K_x=0$ достаточно линейного сравнения.
Каждая принятая проба удовлетворяет
\begin{equation}\label{eq:acceptedprogress}
 f_k-f_{k+1}\ge\alpha_k a_k/2,\qquad
 L_k\le\max\{L_{k-1}/2,2K(\delta_k)\}.
\end{equation}
Значение функции убывает, а нижняя граница не убывает.
Значит, $U$ не возрастает, в том числе при уточнениях в одной точке.

Отнесём обновление к фазе $b_j/2<U_k\le b_j$.
Положительные границы $C_j$ не возрастают.
Если $L_k>C_j$, принятая проба была первой: иначе предыдущее
отклонённое значение было бы не меньше $C_j/2\ge K(\delta_k)$.
Такое обновление уменьшает унаследованную кривизну вдвое.
Перед обновлением фазы $j$ рассмотрим потенциал
\[
 V_k=\max\{0,\log_2(L_{k-1}/C_j)\}.
\]
Обновление выше границы уменьшает его на единицу;
обновление не выше границы оставляет нулевое значение до смены фазы.
Переход из $j$ в $j'$ увеличивает потенциал не более чем на
$\log_2(C_j/C_{j'})$. Эти приращения телескопируются,
поэтому обновлений выше своих границ не больше $P_K$.

На остальных обновлениях \eqref{compact:certified-slope} и
$\|d\|\le D$ дают убывание функции не меньше
\[
 p_j=\min\{s_j/2,s_j^2/(2C_jD^2)\}=s_j^2/(2C_jD^2).
\]
Суммарное убывание от точек фазы $j$ не больше действительной
ошибки при входе в фазу, а она не больше $b_j$.
Таких обновлений не больше $b_j/p_j$.
Суммирование даёт \eqref{nloracle:work}.
Та же конечная граница верна для любого незавершённого префикса,
что исключает бесконечную работу без остановки.

Если на обновлении $k$ отклонено $q_k$ проб, то
$L_k=2^{q_k-1}L_{k-1}$. Следовательно,
\[
 \sum_{k<N}(q_k+1)=2N+\log_2(L_{N-1}/M),\qquad
 L_{N-1}\le\max\{M/2^N,2K(U_0)\}.
\]
Добавим начальное вычисление функции и получим
\eqref{nloracle:functions}.
\end{proof}

\paragraph{Шаги от активных атомов.}\label{alg:afw}
Для $Q=\operatorname{conv}\mathcal A$ с конечным набором атомов
поддерживаем точные положительные веса активного множества $S$.
В AFW полный LMO возвращает минимизирующий атом из $\mathcal A$,
в том числе при неоднозначном минимуме.
Для остановки и проверки точности FW используем $B=h_{\rm FW}+\eta$.
Только после успешной проверки просматриваем активное множество
и сравниваем его приближённый зазор away с зазором FW.
Выбираем больший; при равенстве или одном активном атоме используем FW.
Из выбранного зазора вычитаем $\eta$.
Гарантия ошибки должна одновременно покрывать оба направления,
как в следующем разделе.
Для направления $x-v$ предельный шаг равен $\lambda_v/(1-\lambda_v)$;
заменим им ограничение $1$ в \eqref{eq:btstep}~\cite{lacostejaggi}.
Шаг, достигший этого предела, удаляет атом.
Каждый другой шаг имеет ту же границу продвижения, что и FW.
Потенциальный аргумент считает все обновления выше границ $C_j$,
а аргумент через убывание считает остальные хорошие обновления.
Поэтому хороших обновлений не больше $T_\varepsilon$.
FW добавляет не больше одного атома, а удаление убирает один, откуда
\begin{equation}\label{eq:dropcount}
 N\le2T_\varepsilon+|S_0|-1.
\end{equation}
Границы запросов и значений функции используют фактическое полное $N$.
Дополнительно AFW выполняет $N$ просмотров активного множества.

\begin{corollary}[Скалярное вычисление с доказанной точностью]
\label{nloracle:energy}
Пусть $\delta\le\Phi_0(G)$, а скалярная процедура возвращает
$\Phi_0(B)\le\widehat\Phi_0(B)\le\vartheta\Phi_0(B)$
с известным $\vartheta\ge1$.
Возьмём $C=\widehat\Phi_0(B)$ либо минимум этого значения и
другой корректной верхней границы ошибки в текущей точке.
Теорема применима с $\Phi=\vartheta\Phi_0$.
В частности, можно взять $\Phi_0=\mathcal B$ или $\Phi_0=\mathcal B_D$, а при
$\nabla f(x^*)=0$ также $\Phi_0=W$.
Обратная функция нужна в анализе, а не в алгоритме.
\end{corollary}
\begin{proof}
Корректность даёт $\delta\le C$, а построение даёт
$U\le C\le\vartheta\Phi_0(B)$.
Это именно калибровочные неравенства, использованные в доказательстве.
Нужные свойства $\mathcal B$ и $W$ уже доказаны.
\end{proof}

Нормовая граница $P^+$ позволяет выбрать точный
$C=\mathcal J_D(B,P^+)$.
При вычисленном верхнем включении используем
$C=\widehat{\mathcal J}_D(B,P^+)\le\vartheta\mathcal J_D(B,P^+)
\le\vartheta\mathcal B_D(B)$ и эффективную калибровку
$\Phi=\vartheta\mathcal B_D$.
Последнее неравенство следует из точности $\mathcal J_D$
в классе с верхними границами наблюдений.
Оба варианта можно объединить с другими корректными границами
в текущей точке, взяв их минимум.
Вычисление нормы и скалярного сертификата оплачивается отдельно.
Если на запросе $j$ процедура использует $n(B_j)$ сравнений,
учитывается $\sum_{j=1}^{N+1+R}n(B_j)$, включая неудачные уточнения
и последнюю проверку. Счёт сравнений не ограничивает стоимость
интервальной или битовой арифметики.
Для невозрастающей границы стоимости запроса $C_{\rm query}(\eta)$
общая стоимость не больше
\[
 (N+1)C_{\rm query}(\eta_R)+\sum_{r=0}^{R-1}C_{\rm query}(\eta_0 2^{-r}).
\]

\paragraph{Линейная граница зазора и последние фазы.}
Из известной оценки $G\ge\gamma\delta$, $\gamma>0$, получаем
$\Phi(g)=g/\gamma$.
При $K(u)=H_0+H_1u$ и $\bar\gamma=\kappa\gamma$
\eqref{nloracle:work} даёт
\[
 N\le P_K+\left\lceil
 \frac{4J}{\bar\gamma}+
 \frac{32D^2H_0}{\bar\gamma^2\varepsilon}+
 \frac{16D^2H_1J}{\bar\gamma^2}\right\rceil.
\]
Действительно, $s_j=\bar\gamma b_j/2$ и
$\sum_{j<J}b_j^{-1}<2/\varepsilon$.
Квазар-выпуклость с центром $x^*$
даёт эту оценку через
$\langle\nabla f(x),x-x^*\rangle\ge\gamma\delta(x)$.
Она не требует выпуклости между произвольными парами точек.
Для точной калибровки при сдвинутом RSI фазы с $b_j/2\le g_c$
имеют $\mathcal B^{-1}(b_j/2)=b_j/2$, если положительный порог
$g_c$ существует. Их граница продвижения имеет ту же форму с $\gamma=1$.
Для $\Phi=\vartheta\mathcal B$ соответствующее условие имеет вид
$b_j/(2\vartheta)\le g_c$, а $s_j=\kappa b_j/(2\vartheta)$.
Те же формулы последних фаз верны для $\mathcal B_D$,
поскольку её диапазон обычного зазора совпадает с диапазоном $\mathcal B$.
